\section{Entry from spatially varying metrics}\label{sec:entry}

We specify an open class of initial metrics without prescribing their
finite initial tuple. The tuple is selected by auxiliary flows on
the fixed closed normalization components at $\alpha=0$. Their
use during a short preparation interval will be justified on the
actual smoothing.

\subsection{The selected constant metric on a fixed component}
Let $Y$ be a fixed connected closed Riemann surface of area $A$.
For matrix functions write $\langle f\rangle=A^{-1}\int_Y f\,dA$,
and let $\lambda_1>0$ be the first positive scalar eigenvalue.
Consider the trivial holomorphic bundle in a fixed global flat frame.

\begin{lemma}\label{lem:component-phase}
Let $h(s)$ be HYM flow from a smooth determinant-one metric with
$cI\le h(0)\le C_0I$. There is a unique constant determinant-one
matrix $P$ such that $h(s)\to P$ uniformly. With
$M(s)=\langle h(s)\rangle$, $J(s)=\log\det M(s)$ and
$W(s)=\langle\|h(s)-M(s)\|^2\rangle$, one has
\begin{equation}\label{eq:phase-estimates}
 \begin{gathered}
 M'(s)\preceq0,\qquad P=\lim_{s\to\infty}M(s),\qquad
 cI\le P\le C_0I,\\
 \frac{W}{2C_0^2}\le J\le\frac{W}{2c^2},\qquad
 J(s)\le J(0)e^{-\gamma s},\qquad
 \tr(M(s)-P)\le C_0J(s),\qquad
 \gamma=\frac{2c^2}{C_0^2}\lambda_1.
 \end{gathered}
\end{equation}
For $s\ge1$,
\begin{equation}\label{eq:phase-uniform}
 \sup_Y d(h(s),P)^2\le
 C_Y(2J(0)+J(0)^2)e^{-\gamma(s-1)}.
\end{equation}
If $P,P'$ are selected by two initial metrics in this bounded range,
then
\begin{equation}\label{eq:phase-contraction}
 A D_{\tr}(P,P')\le
          \int_Y D_{\tr}(h(0),h'(0))\,dA.
\end{equation}
Constants in these conclusions do not use derivatives of the
initial metrics.
\end{lemma}

\begin{proof}
Equation \eqref{eq:local-hym} is $h_s=-\Delta h-Q(h)$ with
$Q=4a^{-1}h_z^*h^{-1}h_z\succeq0$. Averaging gives
$M'=-\langle Q\rangle\preceq0$. The bounded decreasing
Hermitian family $M(s)$ has a limit $P$ with eigenvalues in
$[c,C_0]$; the determinant of this limit has not yet been
determined.

For $f(B)=\log\det B$ and a Hermitian matrix $Z$,
\[
 -D^2f_B[Z,Z]=\tr(B^{-1}ZB^{-1}Z),\qquad
 C_0^{-2}\|Z\|^2\le -D^2f_B[Z,Z]\le c^{-2}\|Z\|^2
\]
on the same matrix interval. Taylor's integral formula along the
segment from $M$ to $h(x)$, averaged over $Y$, has zero linear
term. Since $\log\det h(x)=0$, it gives the two bounds between
$J$ and $W$ in \eqref{eq:phase-estimates}.

Using $M^{-1}\ge C_0^{-1}I$, then $h^{-1}\ge C_0^{-1}I$,
gives
\[
 -J'=\tr(M^{-1}\langle Q\rangle)
       \ge C_0^{-1}\langle\tr Q\rangle
       \ge C_0^{-2}\langle|dh|^2\rangle.
\]
Here $4a^{-1}\tr(h_z^*h_z)=|dh|^2$. This identity takes the
trace before any nonconstant matrix weight is introduced; no
weighted commutator has been dropped. Applying the scalar
Poincar\'e inequality to the real matrix entries of $h-M$
gives $\langle|dh|^2\rangle\ge\lambda_1 W\ge2c^2\lambda_1J$.
This proves the exponential bound for $J$.

It follows that $\det P=1$. Since $M-P\succeq0$,
integration of the logarithmic determinant along the segment
from $P$ to $M$ gives
$J\ge C_0^{-1}\tr(M-P)$. The cross term vanishes when averaging
$\|h-P\|^2$, so
\[
 \langle\|h-P\|^2\rangle
    =W+\|M-P\|^2\le C_0^2(2J+J^2).
\]
This proves exponential $L^2$ convergence.

The constant $P$ is an exact stationary solution. By
Lemma~\ref{lem:forced-trace}, $D_{\tr}(h,P)$ is a scalar heat
subsolution. The time-one scalar heat kernel on this fixed
compact surface is bounded. Thus its supremum at $s\ge1$ is at
most a fixed multiple of its integral at $s-1$. Equivalence of
trace and squared matrix distances on the bounded range proves
\eqref{eq:phase-uniform}. Applying the integrated unforced trace
comparison to two flows and passing to their uniform limits
proves \eqref{eq:phase-contraction}.
\end{proof}

The selected matrix depends on the full initial metric. For
simultaneously diagonal data
$h(0)=\operatorname{diag}(e^{f_1},\ldots,e^{f_r})$ with
$\sum_a f_a=0$, the scalar heat equation gives
$P=\operatorname{diag}(e^{\langle f_1\rangle},\ldots,
e^{\langle f_r\rangle})$. In general neither the initial
arithmetic mean nor the exponential of its mean logarithm is
asserted to equal $P$.

\subsection{An open class on the smooth fibres}
Fix $0<\beta<1$. Let $\Omega_i(2\rho)$ be the closed component
with the branch disks of radius $2\rho$ removed, identified with
its region in the smooth fibre. For a smooth trace-free
$h_0$-Hermitian logarithm $S$, express its restriction $S_i$
in the harmonic component gauge and set
\begin{equation}\label{eq:extension-norm}
 \cN(S)=\|S\|_{C^0(C_\tau)}
  +\sum_i\inf_{\widetilde S_i|_{\Omega_i(2\rho)}=S_i}
                    \|\widetilde S_i\|_{C^{6,\beta}(C_i)}.
\end{equation}
Each infimum is over smooth Hermitian trace-free extensions
to the closed component $C_i$.
Define $\mathcal B_{R_0}=\{e^S:\cN(S)<R_0\}$. For each
fixed fibre this is an open nonempty class in the smooth
$C^{6,\beta}$ topology. Indeed restriction to the punctured
components has a bounded linear extension across their smooth
boundary circles, so each infimum seminorm is continuous.
No uniform estimate for this extension operator as
$\rho\to0$ is needed or claimed.

The class includes small smooth component profiles. It also
permits arbitrary smooth modifications inside the removed disks
and middle nodes, provided the global $C^0$ bound and strict
norm margin are preserved. In particular their spatial
derivatives there need not be uniformly bounded. Choose $R_0$
small enough that phases obtained from extensions with total
norm below $2R_0$ lie in the finite ball of
Theorem~\ref{thm:all-time}. This is possible because their
order bounds imply
$d_A(P,o)\le2\sqrt{r\sum_i A_i}\,R_0$.

\begin{theorem}\label{thm:spatial-entry}
For the fixed data of Theorem~\ref{thm:all-time} and sufficiently
small $R_0>0$, every initial metric in $\mathcal B_{R_0}$ has
a selected tuple $P$ and a common preparation time
$s_*\to\infty$ such that $\sigma s_*\to0$ and
\begin{equation}\label{eq:entry-trace}
 \eta_*:=\int_{C_\tau}
             D_{\tr}(h(s_*),\widehat K(P))\,dA_\tau\longrightarrow0.
\end{equation}
The time and convergence bounds are uniform over this class and
all permitted scale ratios. If $H_f(s)$ is the finite solution
with $H_f(0)=P$, then
\begin{equation}\label{eq:spatial-all-time}
 \sup_{s\ge s_*}\int_{C_\tau}
      d(h(s),\widehat K(H_f(s)))^2\,dA_\tau
 \le C\{\eta_*+a_*+(\sigma s_*)^2\}\longrightarrow0.
\end{equation}
Different admissible extensions in \eqref{eq:extension-norm}
give tuples at distance $O(\rho)$ and reconstructions at
integrated squared distance $O(\rho^2)$ at every later time.
\end{theorem}

\begin{proof}
Choose extensions $\widetilde S_i$ with total $C^{6,\beta}$
norm below $2R_0$. Let $\Psi_i(s)$ be HYM flow for the trivial
connection on the fixed closed component, from
$\exp(\widetilde S_i)$, and let $P_i$ be its phase from
Lemma~\ref{lem:component-phase}. These auxiliary flows have
$\alpha_i=0$. Their difference from the actual component
equation will be estimated explicitly.

For each integer $n\ge1$, fixed-component parabolic theory
gives a constant $C_n\ge1$ controlling two spatial derivatives
and one time derivative of all these flows on $[0,n]$.
One can choose $C_n$ nondecreasing. To justify uniformity over
the profiles, the bounded initial $C^{6,\beta}$ class is
precompact in $C^{4,\beta/2}$; continuous dependence for the
fixed strictly parabolic equation, the order bounds and
fixed-component continuation give a finite bound on each
finite time interval. No assertion that $C_n$ stays bounded
as $n\to\infty$ is used.

Construct a preparation metric $K_{\mathrm{pre}}(s)$ with
$\Psi_i(s)$ on $\Omega_i(2\rho)$. Use node geodesics between
their values at the branch points, with parameter
$-\epsilon_e\log|u_e|$, and interpolate relative logarithms
on the incidence collars, including the actual unitary
changes of frame. As before this is a global smooth
determinant-normalized metric with fixed order bounds.
On the component exterior, subtracting the trivial-connection
equation from the actual $D_i$ equation leaves integrated
error at most $C_n\alpha_*$. Metric transfer contributes
$C_n b^2\rho^{-N}$. The node curvature is zero and its time
residual has integral at most $C_n\rho^2$.

On an incidence collar the component value differs from its
value at the branch point by $O(C_n\rho)$, including the
change of frame. The node logarithm differs from that point
value by $O(C_n\epsilon_e|\log\rho|)$. The discrepancy
derivatives of order $k\le2$ are therefore bounded by
$C_n(\rho+\epsilon_*|\log\rho|)\rho^{-k}$. The full
curvature-difference estimate used in Section~\ref{sec:angular}
gives the integrated bound
$C_n(\rho+\epsilon_*|\log\rho|)$, with the quadratic terms
included. Consequently
\begin{equation}\label{eq:preparation-residual}
 \sup_{0\le s\le n}\int_{C_\tau}|\cR(K_{\mathrm{pre}}(s))|
       \le C_n\delta_*,
 \qquad
 \delta_*=\alpha_*+\epsilon_*\ell+b^2\rho^{-N}\to0.
\end{equation}
At time zero $K_{\mathrm{pre}}$ and $h(0)$ agree on every
$\Omega_i(2\rho)$ and have bounded difference elsewhere.
The latter region has area $O(\rho^2)$. Integrated forced
trace comparison gives error at most
$C\rho^2+C_nn\delta_*$ through time $n$.

For sufficiently small parameters choose
\begin{equation}\label{eq:preparation-time}
 s_*=\max\{n\in\mathbb N:nC_n\le\delta_*^{-1/2}\}.
\end{equation}
Then $s_*\to\infty$ and $\sigma s_*\to0$, since
$\sigma\le C\delta_*$. Uniform convergence of the auxiliary
flows to $P_i$, the component Green $L^2$ bounds, and the
area $O(\rho^2)$ of the remaining regions show
\begin{equation}\label{eq:quantitative-entry}
 \eta_*\le C\bigl(\sqrt{\delta_*}
                +e^{-\gamma_0(s_*-1)}+\epsilon_*^2\bigr)
                   \longrightarrow0
\end{equation}
for a fixed $\gamma_0>0$ depending on the component
geometries and initial order bounds. On this range trace
distance and squared affine distance are equivalent, so the
triangle comparison used here costs only a fixed constant.

Restart Theorem~\ref{thm:all-time} at $s_*$ from the finite
tuple $P$. This first compares $h(s)$ with
$\widehat K(H_f(s-s_*))$. Convex gradient flow on the
Hadamard product is nonexpanding: first variation of the
squared distance between two solutions and convexity make
its derivative nonpositive. Autonomy therefore gives
\[
 d_A(H_f(s-s_*),H_f(s))
       \le d_A(P,H_f(s_*))\le C\sigma s_*
       \qquad(s\ge s_*).
\]
Lemma~\ref{lem:reconstruction-energy} transfers this
difference to global integrated squared distance, proving
\eqref{eq:spatial-all-time}.

Two choices of extensions agree on $\Omega_i(2\rho)$.
Their initial component trace difference is bounded and
supported on disks of total area $O(\rho^2)$.
Equation \eqref{eq:phase-contraction} yields
$d_A(P,P')=O(\rho)$. Nonexpansion of finite flow and
\eqref{eq:reconstruction-lipschitz} preserve the asserted
ambiguity for all later times.
\end{proof}

The preparation time is negligible compared with every finite
scale: $\sigma s_*\to0$, while all finite speeds are bounded
by a constant times $\sigma$. No approximation is asserted
during the discarded initial layer. The tuple $P$ depends on
the spatial initial data; convergence to a single
parameter-independent finite trajectory would require
convergence of the selected phases and of the rescaled
parameters, in addition to the estimates proved here.
